% Part: incompleteness
% Chapter: theories-computability
% Section: oconsis-ext-of-q-undec

\documentclass[../../../include/open-logic-section]{subfiles}

\begin{document}

\olfileid[pt]{inc}{tcp}{oqn}

\olsection{Extensões $\omega$-consistentes de $\Th{Q}$ são indecidíveis}

\begin{explain}
A prova de que $\Th{Q}$ é c.e.-completa baseou-se no fato de que toda
sentença demonstrável em $\Th{Q}$ é ``verdadeira'' acerca dos números
naturais. A definição e o teorema a seguir fortalecem esse resultado,
identificando precisamente os aspectos da ``verdade'' necessários na
prova acima. Não se detenha muito nesse teorema, porém, pois em breve o
fortaleceremos ainda mais. Nós o incluímos principalmente por razões
históricas: o artigo original de G\"odel usou a noção de
$\omega$-consistência, mas pouco depois seu resultado foi fortalecido
substituindo a $\omega$-consistência pela consistência ordinária.
\end{explain}

\begin{defn}
\ollabel{thm:oconsis-q}
Uma teoria $\Th{T}$ é $\omega$-consistente se vale o seguinte: se
$\lexists[x][!A(x)]$ é uma sentença qualquer e $\Th{T}$ prova
$\lnot !A(\num 0)$, $\lnot !A(\num 1)$, $\lnot !A(\num 2)$, \dots,
então $\Th{T}$ não prova $\lexists[x][!A(x)]$.
\end{defn}

\begin{thm}
  Seja $\Th{T}$ uma teoria $\omega$-consistente qualquer que inclua
  $\Th{Q}$. Então $\Th{T}$ não é decidível.
\end{thm}

\begin{proof}
Se $\Th{T}$ inclui $\Th{Q}$, então $\Th{T}$ representa as funções e
relações computáveis. Basta modificar a prova anterior. Como acima, se
$x \in K$, então $\Th{T}$ prova $\lexists[s][!A_T(\num
  x,\num x, s)]$. Reciprocamente, suponha que $\Th{T}$ prove
$\lexists[s][!A_T(\num x, \num x, s)]$. Então $x$ deve pertencer a
$K$: caso contrário, não há computação com parada da máquina $x$ com
entrada $x$; como $!A_T$ representa a relação $T$ de Kleene, $\Th{T}$
prova $\lnot !A_T(\num x, \num x, \num 0)$,
$\lnot !A_T(\num x, \num x, \num 1)$, \dots, o que torna $\Th{T}$
$\omega$-inconsistente.
\end{proof}

\end{document}
